\documentclass{article}

\usepackage[utf8]{inputenc}
\usepackage[T1]{fontenc}
\usepackage{lmodern}
\usepackage[french]{babel}
\usepackage{userpackages}
\usepackage{scratch3}
\usepackage{amsmath}
\usepackage{geometry}

\newgeometry{
    lmargin=50pt,
    hmargin=50pt,
    top=50pt,
    textheight=650pt
}

\def\source#1{
    \begin{flushright}
        \itshape\scriptsize Source Brevet #1
    \end{flushright}
}
\def\adapte#1{
    \begin{flushright}
        \itshape\scriptsize Adapté du DNB #1
    \end{flushright}
}

\newcounter{exercice}

\def\newexercice{
    \newpage
    \stepcounter{exercice}
    \begin{flushleft}
        \Large\bfseries Exercice \arabic{exercice}
    \end{flushleft}
}

\begin{document}
\newexercice
    Voici un programme de calcul.
    \begin{center}
        \begin{minipage}{.65\linewidth}
            \begin{simplebox}
                Choisir un nombre.
                \vskip8pt
                Prendre le double du nombre de départ.
                \vskip8pt
                Soustraire 6 au résultat.
                \vskip8pt
                Multiplier le résultat par le nombre de départ.
            \end{simplebox}
        \end{minipage}
    \end{center}\vskip10pt
    \begin{enumerate}[1), font=\bfseries]
        \item Vérifier que si on choisit 4 comme nombre de départ, on obtient
            8.
        \item Appliquer ce programme de calcul au nombre $-3$.
        \item Vous trouverez ci-dessous un script, écrit avec scratch.
        \begin{center}
            \begin{scratch}[corner=1ex, num blocks=true]
                \blockinit{Quand \greenflag est cliqué}
                \blocksensing{Demander \ovalnum{Choisi un nombre} et attendre}
                \blockvariable{mettre \selectmenu{x} à \ovalsensing{réponse}}
                \blockvariable{mettre \selectmenu{y} à \ovaloperator{2*\ovalvariable{\,x\,}}}
                \blockvariable{mettre \selectmenu{z} à \ovaloperator{\ovalvariable{\,y\,}-\ovaloperator{\ovalnum{\phantom{---}}}}}
                \blockvariable{mettre \selectmenu{Résultat} à \ovaloperator{\ovalnum{\phantom{---}}*\ovalnum{\phantom{---}}}}
                \blocklook{dire \ovaloperator{\ovalnum{Le nombre final est } et \ovalvariable{Résultat}} pendant \ovalnum{2} secondes}
            \end{scratch}
        \end{center}
        Compléter les lignes 5 et 6 pour que ce script corresponde au programme
        de calcul.
        \item On veut déterminer le nombre à choisir au départ pour obtenir zéro comme résultat.
        \begin{enumerate}[a), font=\bfseries]
            \item On appelle $x$ le nombre de départ. Exprimer en fonction de
            $x$ le résultat final.
            \item Quel(s) nombre(s) doit-on choisir au départ pour obtenir le
            nombre 0 à l'arrivée ?
        \end{enumerate}
    \end{enumerate}
    \adapte{Métropole 2020}

    \newexercice
    \begin{minipage}[c]{.60\linewidth}
        On cherche à dessiner une éolienne avec le logiciel Scratch ; elle est formée de $3$ pales qui tournent autour d'un axe central.
    \end{minipage}\hfill
    \begin{minipage}[c]{.35\linewidth}
        \begin{center}
            \begin{tikzpicture}[scale=0.6,rotate=90]
                \foreach \i in {0,120,240}{
                    \draw(\i:4)--(\i+20:1)--(\i-20:1)--cycle;
                    \draw(\i:0.95)--(0,0);
                }
            \end{tikzpicture}\par 
        Eolienne modélisée par Scratch
        \end{center}
    \end{minipage}
    ~\par 
    \begin{minipage}[c]{.60\linewidth}
        \begin{enumerate}[1), font=\bfseries]
            \item La figure ci-dessous représente une pale d'éolienne.
            \begin{center}
                \begin{tikzpicture}[scale=0.05]
                    \coordinate[label=left:$A$](A)at(0,0);
                    \coordinate[label=below:$C$](C)at(30,-13);
                    \coordinate[label=right:$B$](B)at(30,0);
                    \coordinate[label=above:$E$](E)at(30,13);
                    \coordinate[label=right:$D$](D)at(179.43,0);
                    \tkzMarkRightAngle[red,size=7](E,B,A)
                    \draw(A)--(B)--(E)--(D)--(C)--(B);
                    \tkzMarkAngle[size=7,mark=none](D,C,B)
                    \tkzLabelAngle[pos=20,below=0.15](D,C,B){$85^\circ$}
                    \tkzMarkSegments[mark=|](E,D D,C)
                    \tkzMarkSegments[mark=s||](E,B B,C)
                    \tkzLabelSegment[sloped](E,D){$150$}
                    \tkzLabelSegment[below,sloped](B,C){$13$}
                    \tkzLabelSegment[sloped](A,B){$30$}
                \end{tikzpicture}
            \end{center}
            \begin{itemize}
                \item $DEC$ est un triangle isocèle en $D$~;
                \item $B$ est le milieu de $[EC]$~;
                \item $[AB]$ est perpendiculaire à $[EC]$~; 
                \item $\widehat{ECD}= 85^\circ$.
            \end{itemize}
            \begin{enumerate}[a), font=\bfseries]
                \item Montrer que l'angle $\widehat{\text{CDE}} = 10\degres$.
                \item Le script \og pale \fg{} ci-contre permet de tracer une pale de l'éolienne avec le logiciel Scratch.\par                     
                Pourquoi la valeur indiquée dans le bloc de la ligne \no 6 est-elle 95~?
                \item Dans ce même script \og pale \fg{}, par quelle valeur doit-on compléter le bloc situé à la ligne \no 8~? \par                     
                Recopier cette valeur sur votre copie.
            \end{enumerate}
            \item Le script \og éolienne \fg{} ci-contre permet de tracer l'éolienne avec le logiciel Scratch. \par                 
            Par quelle valeur doit-on compléter la boucle \og  répéter\fg{}~? Recopier cette valeur sur votre copie.                
        \end{enumerate}
    \end{minipage}\hfill
    \begin{minipage}[c]{.35\linewidth}
        \begin{center}
            \setscratch{scale=.75}
            \begin{scratch}[num blocks]
                \initmoreblocks{définir \namemoreblocks{pale}}
                \blockpen{stylo en position écriture}
                \blockmove{avancer de \ovalnum{30}}
                \blockmove{tourner \turnright{} de \ovalnum{90} degrés}
                \blockmove{avancer de \ovalnum{13}}
                \blockmove{tourner \turnleft{} de \ovalnum{95} degrés}
                \blockmove{avancer de \ovalnum{150}}
                \blockmove{tourner \turnleft{} de \ovalnum{} degrés}
                \blockmove{avancer de \ovalnum{150}}
                \blockmove{tourner \turnleft{} de \ovalnum{95} degrés}
                \blockmove{avancer de \ovalnum{13}}
                \blockmove{tourner \turnright{} de \ovalnum{90} degrés}
                \blockmove{avancer de \ovalnum{30}}
                \blockmove{tourner \turnright{} de \ovalnum{180} degrés}
                \blockpen{relever le stylo}
            \end{scratch}
            ~\par \medskip
            \begin{scratch}
                \initmoreblocks{définir \namemoreblocks{éolienne}}
                \blockmove{aller à x: \ovalnum0 y: \ovalnum0}
                \blockrepeat{répéter \ovalnum{} fois}
                {\blockmoreblocks{pale}
                    \blockmove{tourner \turnright{} de \ovalnum{120} degrés}
                }
            \end{scratch}
        \end{center}
    \end{minipage}
    \source{Métropole 2020}
\end{document}